\section{Exact certificates for the remaining golden weight-five ladders}
\label{gold:sec}

Put
\[
 \rho=\frac{\sqrt5-1}{2},\qquad \phi=\rho^{-1},\qquad
 \lambda=\log\rho=-\log\phi.
\]
The golden-ratio discussion in the audited manuscript identifies three
weight-five evaluations with the index-$12$, $20$, and $24$ ladder
families, and records numerical evidence for them \cite{Repo}.
The first evaluation has already received an analytic proof in the
preceding research.  The contribution here is a complete local proof
certificate for each of the other two evaluations, together with all
its lower-weight companions.  These are statements about what is now
proved explicitly in the repository's research program; they do not
assert historical priority for the ladder evaluations.  The general
connection between rational-function identities, tensor kernels, and
polylogarithm ladders is classical; see Gangl \cite{Gangl}.

The certificates use rational functions of one variable.  Consequently,
the lower-weight conclusions follow from identities holding on an open
interval.  No differentiation of an identity valid only at the isolated
algebraic number $\rho$ is involved.

\subsection{The two evaluations and their full descent}

Define two finitely supported coefficient vectors by
\begin{align}
 (c^{\mathrm{II}}_2,c^{\mathrm{II}}_4,c^{\mathrm{II}}_{10},
     c^{\mathrm{II}}_{20})&=(45000,-16875,-144,9),
 \label{gold:eq:cII}\\
 (c^{\mathrm{III}}_2,c^{\mathrm{III}}_4,c^{\mathrm{III}}_6,
     c^{\mathrm{III}}_8,c^{\mathrm{III}}_{24})
     &=(15660,-19440,7680,2430,-15).
 \label{gold:eq:cIII}
\end{align}
A missing coefficient means zero.  The superscripts refer to the
second and third displayed golden weight-five formulas in the
manuscript; the third vector is the displayed formula before the
canonical index-$24$ recombination.

\begin{theorem}[Golden evaluations with exact descent]
\label{gold:thm:main}
For either vector $c=c^{\mathrm{II}},c^{\mathrm{III}}$, put
\[
 S_m=\sum_{j\ge1}c_jj^{5-m}\Li_m(\rho^j),\qquad 1\le m\le5.
\]
Then
\begin{equation}
 S_m=\frac{a\lambda^m}{m!}
     +\sum_{r=2}^{m}q_r\zeta(r)\frac{\lambda^{m-r}}{(m-r)!},
 \label{gold:eq:all-weights}
\end{equation}
where the constants are
\begin{center}
\begin{tabular}{c|rrrrr}
 & $a$ & $q_2$ & $q_3$ & $q_4$ & $q_5$\\ \hline
 $\mathrm{II}$ & $720000$ & $-57600$ & $0$ & $32040$ & $28944$\\
 $\mathrm{III}$ & $-250560$ & $8640$ & $0$ & $2520$ & $5025$
\end{tabular}
\end{center}
In particular, writing $\ell=\log\phi$, the two weight-five evaluations are
\begin{align}
 &45000\Li_5(\rho^2)-16875\Li_5(\rho^4)
       -144\Li_5(\rho^{10})+9\Li_5(\rho^{20})\notag\\
 &\hspace{12mm}=28944\zeta(5)-6000\ell^5
           +1600\pi^2\ell^3-356\pi^4\ell,
 \label{gold:eq:II-five}\\[2mm]
 &15660\Li_5(\rho^2)-19440\Li_5(\rho^4)
       +7680\Li_5(\rho^6)+2430\Li_5(\rho^8)
       -15\Li_5(\rho^{24})\notag\\
 &\hspace{12mm}=5025\zeta(5)+2088\ell^5
          -240\pi^2\ell^3-28\pi^4\ell.
 \label{gold:eq:III-five}
\end{align}
\end{theorem}

The vanishing $q_3=0$ matters: the descent produces no
$\zeta(3)\lambda^2$ term at weight five.  This cancellation will follow
from the exact endpoint arithmetic of the functional certificates.

\subsection{A real generalized Rogers function}

For $n\ge2$ and real $x\ne0,1$, define
\begin{align}
 \mathcal R_n(x)
 &=\sum_{r=0}^{n-2}\frac{(-\log|x|)^r}{r!}
                  \operatorname{Re}\Li_{n-r}(x)\notag\\
 &\quad+\frac{(-1)^{n-1}(n-1)}{n!}
           \log^{n-1}|x|\operatorname{Re}\Li_1(x).
 \label{gold:eq:R-definition}
\end{align}
Here the ordinary polylogarithms use their principal branches, and
$\operatorname{Re}\Li_1(x)=-\log|1-x|$.  We extend the definition
continuously by $\mathcal R_n(0)=0$ and $\mathcal R_n(1)=\zeta(n)$.
The first functions are
\begin{align*}
 \mathcal R_2(x)&=\operatorname{Re}\Li_2(x)
                +\tfrac12\log|x|\log|1-x|,\\
 \mathcal R_3(x)&=\operatorname{Re}\Li_3(x)
                -\log|x|\operatorname{Re}\Li_2(x)
                -\tfrac13\log^2|x|\log|1-x|,\\
 \mathcal R_4(x)&=\operatorname{Re}\Li_4(x)
                -\log|x|\operatorname{Re}\Li_3(x)
                +\tfrac12\log^2|x|\operatorname{Re}\Li_2(x)
                -\tfrac18\log^3|x|\operatorname{Re}\Li_1(x),\\
 \mathcal R_5(x)&=\operatorname{Re}\Li_5(x)
                -\log|x|\operatorname{Re}\Li_4(x)
                +\tfrac12\log^2|x|\operatorname{Re}\Li_3(x)\\
 &\qquad-\tfrac16\log^3|x|\operatorname{Re}\Li_2(x)
                +\tfrac1{30}\log^4|x|\operatorname{Re}\Li_1(x).
\end{align*}
In particular, $\mathcal R_3$ is the usual real single-valued
trilogarithm.  At higher weights it is useful to retain
$\mathcal R_n$ itself, including at even weights.

\begin{lemma}[Derivative, inversion, and duplication]
\label{gold:lem:R}
Write $L=\log|x|$ and $M=\log|1-x|$.  On each real interval avoiding
$0$ and $1$,
\begin{equation}
 d\mathcal R_n(x)
   =\frac{(-1)^n(n-1)}{n!}
       L^{n-2}\bigl(L\,dM-M\,dL\bigr).
 \label{gold:eq:R-derivative}
\end{equation}
Let
\[
 h_n(+1)=1,\qquad h_n(-1)=-(1-2^{1-n}).
\]
For nonzero real $x$,
\begin{align}
 \mathcal R_n(x^{-1})
 &=(-1)^{n-1}\mathcal R_n(x)
   +\bigl(1-(-1)^{n-1}\bigr)h_n(\operatorname{sgn}x)\zeta(n),
 \label{gold:eq:R-inversion}\\
 \mathcal R_n(x)+\mathcal R_n(-x)
 &=2^{1-n}\mathcal R_n(x^2).
 \label{gold:eq:R-duplication}
\end{align}
\end{lemma}

\begin{proof}
The derivative relation
$d\operatorname{Re}\Li_j(x)=\operatorname{Re}\Li_{j-1}(x)\,dL$
holds for $j\ge2$.  Consequently the derivative of
\[
 F_n(x)=\sum_{r=0}^{n-1}\frac{(-L)^r}{r!}
                      \operatorname{Re}\Li_{n-r}(x)
\]
telescopes to $(-1)^nL^{n-1}dM/(n-1)!$.
Since
\[
 \mathcal R_n=F_n+\frac{(-1)^{n-1}}{n!}L^{n-1}M,
\]
one more differentiation proves \eqref{gold:eq:R-derivative}.
Under inversion, $(L,M)$ becomes $(-L,M-L)$.  Substitution in the
derivative formula shows that the difference in
\eqref{gold:eq:R-inversion} is constant separately on each of the
positive and negative half-axes.  Evaluation at $x=1$ and $x=-1$
fixes the constants, because
\[
 \mathcal R_n(-1)=\Li_n(-1)=-(1-2^{1-n})\zeta(n).
\]
The duplication formula follows term by term from the ordinary
polylogarithm duplication formula and $\log(x^2)=2\log|x|$.
These arguments extend across $x=1$ by continuity.  In
\eqref{gold:eq:R-definition}, the potentially singular monologarithm
is multiplied by a positive power of $\log|x|$, and its product tends
to zero as $x\to1$.
\end{proof}

\subsection{A tensor certificate with an elementary analytic interpretation}

Let $V=\mathbb Q(t)^\times\otimes_{\mathbb Z}\mathbb Q$, with the
multiplicative group written additively.  Signs are torsion and therefore
vanish in $V$.  For a nonzero rational function $f\ne1$, put
\[
 \beta_n[f]=f^{\otimes(n-2)}\otimes\bigl(f\wedge(1-f)\bigr).
\]
The repeated first factors can equivalently be placed in a symmetric
tensor power.  Only elementary linear algebra with this notation is
needed below.

The rational functions in the certificates have the form
\begin{equation}
 f_{s,a,b,c}(t)=s\,t^a(1-t)^b(1+t)^c,
 \qquad s\in\{1,-1\}.
 \label{gold:eq:arguments}
\end{equation}
They have neither zeros nor poles for $0<t<1$.  Define
\begin{equation}
 k(a,b,c)=a+2b-c.
 \label{gold:eq:k}
\end{equation}
This integer is the exponent after golden specialization, since
$1-\rho=\rho^2$ and $1+\rho=\rho^{-1}$.

\begin{lemma}[Exact descent from a rational-function certificate]
\label{gold:lem:descent}
Suppose $f_i=f_{s_i,a_i,b_i,c_i}$ and $u_i\in\mathbb Q$ satisfy
\begin{equation}
 \sum_i u_i\beta_5[f_i]=0.
 \label{gold:eq:tensor-certificate}
\end{equation}
Write $k_i=a_i+2b_i-c_i$.  For $2\le n\le5$, the function
\begin{equation}
 H_n(t)=\sum_i u_i k_i^{5-n}\mathcal R_n(f_i(t))
 \label{gold:eq:functional-family}
\end{equation}
is constant on $0<t<1$.  If all $a_i\ge0$, its value is
\begin{equation}
 H_n(t)=E_n\zeta(n),\qquad
 E_n=\sum_{a_i=0}u_i k_i^{5-n}h_n(s_i).
 \label{gold:eq:E}
\end{equation}
We use $k_i^0=1$, including when $k_i=0$.
\end{lemma}

\begin{proof}
Extend the assignments
\[
 \eta(t)=1,\qquad \eta(1-t)=2,\qquad \eta(1+t)=-1
\]
to a linear functional $\eta:V\to\mathbb Q$.  For example, assign
zero to every other primitive irreducible polynomial and every rational
prime.  Then $\eta(f_i)=k_i$.  Contracting $5-n$ of the first three
tensor factors in \eqref{gold:eq:tensor-certificate} gives
\begin{equation}
 \sum_i u_i k_i^{5-n}\beta_n[f_i]=0.
 \label{gold:eq:contracted}
\end{equation}
For a fixed regular $t$, evaluate a tensor factor $g$ by $\log|g(t)|$,
and evaluate a wedge $g\wedge h$ by
\[
 \log|g(t)|\,d\log|h(t)|
 -\log|h(t)|\,d\log|g(t)|.
\]
Equation \eqref{gold:eq:contracted}, together with
\eqref{gold:eq:R-derivative}, gives $dH_n=0$.
There are only finitely many points in $(0,1)$ where some $f_i(t)=1$.
The functions $\mathcal R_n(f_i(t))$ are continuous at these points,
so the constants on the intervening intervals agree.
Finally, if $a_i>0$ then $f_i(t)\to0$ as $t\to0^+$, while if
$a_i=0$ then $f_i(t)\to s_i$.  These limits prove \eqref{gold:eq:E}.
\end{proof}

A rational constant in $1-f_i$ must be retained when verifying the
tensor equation.  For example,
\[
 1-\frac{1-t}{1+t}=\frac{2t}{1+t}
\]
contributes the independent constant-prime coordinate $2$ to $1-f_i$.
It has logarithm $\log2$ and derivative zero; discarding it would change
both the tensor equation and its analytic interpretation.  The
certificates below and the verification program retain this coordinate.

\subsection{The two finite certificates}

Table~\ref{gold:tab:certificates} defines coefficients $u_{s,a,b,c}$
and $v_{s,a,b,c}$.  An entry $0$ means that the corresponding argument
is absent.  There are $58$ nonzero $u$ entries and $60$ nonzero $v$
entries.  For each column separately,
\begin{equation}
 \sum_{s,a,b,c}u_{s,a,b,c}\beta_5[f_{s,a,b,c}]=0,
 \qquad
 \sum_{s,a,b,c}v_{s,a,b,c}\beta_5[f_{s,a,b,c}]=0.
 \label{gold:eq:two-certificates}
\end{equation}
Thus the table also gives two explicit rational-function identities
\eqref{gold:eq:functional-family}, and four weights of each identity.
Their endpoint constants are displayed below.

Here is a completely finite interpretation of the two checks in
\eqref{gold:eq:two-certificates}.  Factor every $f$ and $1-f$ over
$\mathbb Q(t)$, including the rational contents.  Use rational primes
and primitive irreducible polynomials as a multiplicative basis modulo
signs.  Write the factor vectors as $e_f$ and $d_f$.  For each unordered
pair of basis indices $(p,q)$ and each triple $(i,j,k)$ chosen with
repetition from the three coordinates $t,1-t,1+t$, verify
\begin{equation}
 \sum_f u_f(e_f)_i(e_f)_j(e_f)_k
 \bigl((e_f)_p(d_f)_q-(e_f)_q(d_f)_p\bigr)=0,
 \label{gold:eq:finite-coordinate-check}
\end{equation}
and then the same equation with $u$ replaced by $v$.
The first column has $16$ basis coordinates, including the prime $2$;
the second has $18$, also including $2$.  There are respectively
$330$ and $370$ potentially nonzero weight-five coordinate sums, and
every one is zero.  These are exact integer equalities.
The supplied verifier factors the functions afresh and performs these
sums using rational arithmetic.  It additionally checks all four
contracted tensor equations directly.

\begingroup
\small
\setlength{\tabcolsep}{9pt}
\begin{longtable}{rrrrrr}
\caption{The two weight-five rational-function certificates.  The
argument in a row is $s\,t^a(1-t)^b(1+t)^c$.}
\label{gold:tab:certificates}\\
\hline
$s$&$a$&$b$&$c$&$u$&$v$\\ \hline
\endfirsthead
\multicolumn{6}{c}{Table~\ref{gold:tab:certificates} continued}\\
\hline $s$&$a$&$b$&$c$&$u$&$v$\\ \hline
\endhead
\hline
\endfoot
1&0&0&1&-2166&-1161\\
-1&0&0&1&-3240&-2025\\
-1&0&0&3&40&25\\
1&0&1&-1&-80&-288\\
-1&0&1&-1&-80&-288\\
1&0&1&0&2178&3510\\
-1&0&1&0&2916&4536\\
1&0&1&1&6&-270\\
1&0&1&2&26&72\\
-1&0&2&-2&10&36\\
1&0&2&1&-54&-27\\
-1&0&3&0&-36&-56\\
-1&1&-3&-2&8&9\\
1&1&-2&-3&-4&-9\\
-1&1&-2&-1&-40&0\\
-1&1&-2&0&284&405\\
1&1&-1&-1&-6&27\\
-1&1&-1&-1&-162&-351\\
1&1&-1&0&2268&2997\\
-1&1&-1&0&1722&2430\\
-1&1&-1&1&-105&-360\\
1&1&0&-2&-320&-144\\
1&1&0&-1&-1710&-81\\
-1&1&0&-1&-2592&-486\\
1&1&0&0&-1236&-54\\
-1&1&0&0&528&1539\\
1&1&0&1&-126&-414\\
-1&1&0&1&-480&-405\\
1&1&1&-1&-70&-270\\
1&1&1&0&444&909\\
-1&1&1&0&-42&9\\
1&1&1&1&-72&-81\\
-1&1&1&1&42&81\\
-1&1&2&1&0&9\\
-1&1&3&1&5&0\\
-1&2&-2&-1&-72&-81\\
-1&2&-1&-2&32&81\\
-1&2&-1&-1&6&-270\\
1&2&-1&0&-206&-504\\
-1&2&-1&0&204&234\\
1&2&0&-1&38&99\\
-1&2&0&-1&-240&270\\
-1&2&0&0&-210&-540\\
1&2&1&0&-60&0\\
-1&2&1&0&0&81\\
1&3&-3&-3&0&-1\\
1&3&-3&0&-28&-37\\
1&3&-2&-1&-46&-9\\
-1&3&-1&-2&18&54\\
-1&3&-1&-1&-30&0\\
-1&3&-1&0&-72&-81\\
1&3&-1&1&14&45\\
-1&3&0&-3&32&6\\
1&3&0&-1&12&81\\
-1&3&0&-1&0&81\\
-1&3&0&0&4&19\\
1&4&-3&-2&4&0\\
-1&4&-1&-2&0&9\\
-1&4&0&1&8&9\\
-1&4&1&0&-8&-9\\
1&5&-3&-1&8&9\\
-1&5&-1&-3&-3&-9\\
1&5&1&-1&0&9\\
-1&6&3&0&0&-1\\
-1&7&1&-1&1&0\\
-1&9&0&-3&0&-1\\
\end{longtable}
\endgroup

The table was obtained by exact linear algebra on the functions
\eqref{gold:eq:arguments} with numerator and denominator degrees at most
$9$, after cancellation of common factors.  The search discovers the
rows; the explicit table and the coordinate equations
\eqref{gold:eq:finite-coordinate-check} are the proof certificate.
A complete classification of all functional identities in this or a
larger alphabet is not needed for the argument.

\subsection{Golden specialization and the endpoint constants}

At $t=\rho$, the table argument becomes $s\rho^k$, with
$k=a+2b-c$.  Applying inversion to negative exponents and duplication
to negative arguments gives
\begin{align}
 \sum_f u_f k_f^{5-n}\mathcal R_n(f(\rho))
 &=\frac1{144}\sum_{j\ge1}c^{\mathrm{II}}_j
           j^{5-n}\mathcal R_n(\rho^j)
      +B_n^{\mathrm{II}}\zeta(n),
 \label{gold:eq:u-image}\\
 \sum_f v_f k_f^{5-n}\mathcal R_n(f(\rho))
 &=\frac1{120}\sum_{j\ge1}c^{\mathrm{III}}_j
           j^{5-n}\mathcal R_n(\rho^j)
      +B_n^{\mathrm{III}}\zeta(n).
 \label{gold:eq:v-image}
\end{align}
The constant $B_n$ includes the inversion constants at even weights
and any terms with $k=0$ at weight five.  All remaining coefficients
in \eqref{gold:eq:u-image}--\eqref{gold:eq:v-image} are exact rational
row sums.  The boundary computation is
\begin{center}
\small
\begin{tabular}{c|rrr|rrr}
 &\multicolumn{3}{c|}{$u$ column}&\multicolumn{3}{c}{$v$ column}\\
$n$&$E_n$&$B_n$&$E_n-B_n$&$E_n$&$B_n$&$E_n-B_n$\\ \hline
2&$7122$&$7522$&$-400$&$7695$&$7623$&$72$\\
3&$0$&$0$&$0$&$0$&$0$&$0$\\
4&$-2721/2$&$-1583$&$445/2$&$-7269/4$&$-7353/4$&$21$\\
5&$2205/8$&$597/8$&$201$&$-1011/4$&$-2357/8$&$335/8$
\end{tabular}
\end{center}
For example, the weight-five rational-function identities themselves are
\begin{equation}
 \sum_f u_f\mathcal R_5(f(t))=\frac{2205}{8}\zeta(5),\qquad
 \sum_f v_f\mathcal R_5(f(t))=-\frac{1011}{4}\zeta(5),
 \quad 0<t<1.
 \label{gold:eq:new-functional-identities}
\end{equation}
The golden power with exponent zero contributes $597\zeta(5)/8$ in
the first specialization and $-2357\zeta(5)/8$ in the second.
Subtracting these terms gives $201\zeta(5)$ and
$335\zeta(5)/8$, respectively.  Multiplication by $144$ and $120$
therefore gives the weight-five constants $28944$ and $5025$.

For transparency, the power and sign bookkeeping in every weight can
be expressed without any polylogarithm manipulation.  A row of
coefficient $w$ has weighted coefficient $w k^{5-n}$.  If $k<0$ and
$n$ is even, add $2wk^{5-n}h_n(s)$ to $B_n$ and negate the remaining
coefficient.  Replace $k$ by $|k|$.  If $s=-1$, replace its remaining
coefficient $z$ at $-\rho^k$ by $-z$ at $\rho^k$ and
$2^{1-n}z$ at $\rho^{2k}$.  If $k=0$, put the resulting constant into
$B_n$.  This finite procedure gives every entry above.
Consequently
\begin{equation}
 \sum_{j\ge1}c_jj^{5-n}\mathcal R_n(\rho^j)=q_n\zeta(n),
 \qquad 2\le n\le5,
 \label{gold:eq:golden-R}
\end{equation}
with the constants stated in Theorem~\ref{gold:thm:main}.

\subsection{Removing the lower polylogarithms}

The weight-one companion can be computed using the following elementary
factorizations.  Each equality follows by reduction modulo
$\rho^2+\rho-1=0$, with $\sqrt5=2\rho+1$:
\begin{center}
\begin{tabular}{c|l@{\qquad}c|l}
$j$&$1-\rho^j$&$j$&$1-\rho^j$\\ \hline
$2$&$\rho$&$10$&$11\rho^5$\\
$4$&$\sqrt5\rho^2$&$12$&$8\sqrt5\rho^6$\\
$6$&$4\rho^3$&$20$&$55\sqrt5\rho^{10}$\\
$8$&$3\sqrt5\rho^4$&$24$&$144\sqrt5\rho^{12}$
\end{tabular}
\end{center}
For either vector $c$, the coefficients of $\log2$, $\log3$,
$\log5$, and $\log11$ in
$\sum_j c_jj^4\log(1-\rho^j)$ vanish.  Hence
\begin{equation}
 S_1=-\sum_jc_jj^4\log(1-\rho^j)
     =-\frac12\sum_jc_jj^5\lambda=a\lambda.
 \label{gold:eq:S1}
\end{equation}
The two sums $\sum_j c_jj^5$ are $-1440000$ and $501120$,
which give $a=720000$ and $a=-250560$.

\begin{proof}[Proof of Theorem~\ref{gold:thm:main}]
The finite polynomial identities \eqref{gold:eq:two-certificates},
Lemma~\ref{gold:lem:descent}, and the exact specialization arithmetic
prove \eqref{gold:eq:golden-R}.  Expand the definition of
$\mathcal R_n$ at $\rho^j$.  Since $\log(\rho^j)=j\lambda$, it gives
\begin{equation}
 q_n\zeta(n)=
 \sum_{r=0}^{n-2}\frac{(-\lambda)^r}{r!}S_{n-r}
 +\frac{(-1)^{n-1}(n-1)}{n!}\lambda^{n-1}S_1.
 \label{gold:eq:triangular}
\end{equation}
Starting with \eqref{gold:eq:S1}, this triangular system has the unique
solution \eqref{gold:eq:all-weights}.  One may also verify the formula
by substitution: the coefficients of each $q_j\zeta(j)$ reduce to
\[
 \sum_{r=0}^{n-j}\frac{(-1)^r}{r!(n-j-r)!}
   =\begin{cases}1,&j=n,\\0,&j<n,\end{cases}
\]
and the $a$ term cancels because
\[
 \sum_{r=0}^{n-2}\frac{(-1)^r}{r!(n-r)!}
   +\frac{(-1)^{n-1}(n-1)}{n!}=0.
\]
This proves all five weights.  Substituting
$\zeta(2)=\pi^2/6$, $\zeta(4)=\pi^4/90$, and $\lambda=-\ell$ in
the case $n=5$ proves the two asserted evaluations.
\end{proof}

Some useful lower-weight consequences, written with smaller integer
coefficient vectors, are
\begin{align}
 5\Li_2(\rho^2)-15\Li_2(\rho^4)-2\Li_2(\rho^{10})
    +\Li_2(\rho^{20})
 &=5\lambda^2-\frac{2\pi^2}{15},
 \label{gold:eq:II-two}\\
 50\Li_3(\rho^2)-75\Li_3(\rho^4)-4\Li_3(\rho^{10})
    +\Li_3(\rho^{20})
 &=\frac{100\lambda^3-8\pi^2\lambda}{3},
 \label{gold:eq:II-three}\\
 500\Li_4(\rho^2)-375\Li_4(\rho^4)-8\Li_4(\rho^{10})
    +\Li_4(\rho^{20})
 &=\frac{500\lambda^4}{3}-\frac{80\pi^2\lambda^2}{3}
    +\frac{89\pi^4}{45}.
 \label{gold:eq:II-four}
\end{align}
For the third displayed ladder one obtains
\begin{align}
 &87\Li_2(\rho^2)-864\Li_2(\rho^4)+1152\Li_2(\rho^6)
       +864\Li_2(\rho^8)-144\Li_2(\rho^{24})
    =\pi^2-87\lambda^2,
 \label{gold:eq:III-two}\\
 &87\Li_3(\rho^2)-432\Li_3(\rho^4)+384\Li_3(\rho^6)
       +216\Li_3(\rho^8)-12\Li_3(\rho^{24})\notag\\
 &\hspace{25mm}=-58\lambda^3+2\pi^2\lambda,
 \label{gold:eq:III-three}\\
 &87\Li_4(\rho^2)-216\Li_4(\rho^4)+128\Li_4(\rho^6)
       +54\Li_4(\rho^8)-\Li_4(\rho^{24})\notag\\
 &\hspace{25mm}=-29\lambda^4+2\pi^2\lambda^2+\frac{7\pi^4}{90}.
 \label{gold:eq:III-four}
\end{align}
These companions are exact consequences of the same certificate, with
no separate recognition of numerical constants.

\subsection{The first ladder and canonical normalization}

The first golden formula is included in the supplied replay data as a
small consistency and normalization certificate.  Its ordinary
coefficient vector is
\[
 (c^{\mathrm I}_2,c^{\mathrm I}_4,c^{\mathrm I}_6,c^{\mathrm I}_{12})
       =(4455,-1215,-360,15).
\]
A primitive functional certificate has the following twelve rows in
the notation \eqref{gold:eq:arguments}:
\begin{center}
\small
\begin{tabular}{rrrr|r@{\qquad}rrrr|r}
$s$&$a$&$b$&$c$&$w$&$s$&$a$&$b$&$c$&$w$\\ \hline
$1$&$0$&$0$&$1$&$-54$&$1$&$1$&$0$&$0$&$-81$\\
$-1$&$0$&$0$&$1$&$-81$&$-1$&$1$&$0$&$0$&$-54$\\
$-1$&$0$&$0$&$3$&$1$&$-1$&$1$&$0$&$1$&$-9$\\
$1$&$1$&$0$&$-2$&$-9$&$-1$&$2$&$0$&$-1$&$-9$\\
$1$&$1$&$0$&$-1$&$-54$&$-1$&$3$&$0$&$-3$&$1$\\
$-1$&$1$&$0$&$-1$&$-81$&$1$&$3$&$0$&$0$&$1$
\end{tabular}
\end{center}
Its scale factor is $240$.  The same calculation gives
\[
 (a,q_2,q_3,q_4,q_5)=(84240,-6480,0,3420,3015),
\]
recovering the previously proved first evaluation.

For completeness, define the manuscript's canonical ladders explicitly.
For $N\in\{12,20,24\}$, write
\[
 T_N(m)=D_{N,1}\frac{\lambda^m}{m!}
       +D_{N,2}\zeta(2)\frac{\lambda^{m-2}}{(m-2)!}
       +D_{N,3}\zeta(4)\frac{\lambda^{m-4}}{(m-4)!},
\]
with any term having a negative factorial index omitted, and put
\begin{align*}
 L_{12}(m)&=\frac{\Li_m(\rho^{12})}{12^{m-1}}
 -\frac32\frac{\Li_m(\rho^6)}{6^{m-1}}
 -\frac{\Li_m(\rho^4)}{4^{m-1}}
 +\frac{11}{48}\frac{\Li_m(\rho^2)}{2^{m-1}}+T_{12}(m),\\
 L_{20}(m)&=\frac{\Li_m(\rho^{20})}{20^{m-1}}
 -\frac{\Li_m(\rho^{10})}{10^{m-1}}
 -3\frac{\Li_m(\rho^4)}{4^{m-1}}
 +\frac12\frac{\Li_m(\rho^2)}{2^{m-1}}+T_{20}(m),\\
 L_{24}(m)&=\frac{\Li_m(\rho^{24})}{24^{m-1}}
 -\frac43\frac{\Li_m(\rho^{12})}{12^{m-1}}
 -2\frac{\Li_m(\rho^8)}{8^{m-1}}\\
 &\quad+\frac73\frac{\Li_m(\rho^4)}{4^{m-1}}
 -\frac{205}{576}\frac{\Li_m(\rho^2)}{2^{m-1}}+T_{24}(m).
\end{align*}
The correction coefficients are
\[
\begin{array}{c|rrr}
 N&D_{N,1}&D_{N,2}&D_{N,3}\\ \hline
 12&-13/48&1/48&-19/1728\\
 20&-1/2&1/25&-89/4000\\
 24&179/576&-5/192&629/41472
\end{array}
\]
With these definitions, exact equality holds in
\[
 L_{12}(5)=\frac{67}{6912}\zeta(5),\qquad
 L_{20}(5)=\frac{201}{10000}\zeta(5),\qquad
 L_{24}(5)=-\frac{1541}{110592}\zeta(5).
\]
For the last value, add $64/3$ times the first ordinary weight-five
formula to the third, and divide by $-15\cdot24^4$; this eliminates
$\Li_5(\rho^6)$ and gives the manuscript's canonical index-$24$
combination.  The corresponding normalized first-degree ladders vanish
at weights $1,2,3,4$.  These statements follow by exact rational
recombination of \eqref{gold:eq:all-weights} and its first-ladder
counterpart.

\subsection{Verification, integration, and further questions}

The files \texttt{golden\_rational\_certificates.json} and
\texttt{certify\_golden\_ladders.py} contain the complete finite data and
an independent replay of the polynomial checks.  The verifier uses
SymPy to factor every rational function, retains rational-prime
coordinates, checks the tensor equations directly in all four weights,
performs every golden specialization, and checks the eight elementary
factorizations used for $S_1$.  Its separate $85$-digit calculations
verify the rational-function identities at $t=37/100$ and the ordinary
companions at $\rho$; these numerical diagnostics do not enter the proof.
The detailed receipt is
\texttt{golden\_certificate\_verification.json}.

For manuscript integration, the second and third displayed weight-five
formulas can be marked proved and linked to this section.  Their
first-degree lower-weight vanishing and the canonical values just
displayed can be marked proved as well.  The weight-$6$ through
weight-$9$ reconstructions have a different status: the present
certificates establish no upward continuation beyond weight five.
The contraction in Lemma~\ref{gold:lem:descent} decreases weight;
there is no inverse contraction that would produce the missing higher
functional identity.

Several follow-up questions are concrete.  First, can the $58$- and
$60$-term certificates be decomposed into a short list of named Kummer
specializations, giving a smaller human derivation?  Second, can a
larger rational-function alphabet produce the manuscript's weight-$6$
and weight-$7$ golden identities with similarly complete endpoint
bookkeeping?  Third, which part of the golden relation space is obtained
from a prescribed degree bound, and can exact matrix-rank certificates
prove useful minimality statements within that finite search model?
Fourth, the rational-function identities
\eqref{gold:eq:new-functional-identities} can be specialized at other
algebraic values satisfying multiplicative relations among
$t,1-t,1+t$; determining when their descended constants close on a
small polylogarithm basket is a systematic route to additional ladders.
